La perte individuelle issue du modèle gaussien est

\[
L_i(\theta)
=
L(y_i,Q_\theta(s_i,a_i))
=
\frac{1}{2\widetilde\sigma^2}
\bigl(y_i-Q_\theta(s_i,a_i)\bigr)^2.
\]

La loss agrégée initiale est

\[
\mathcal L(\theta)
=
\frac1n\sum_{i=1}^n L_i(\theta),
\]

tandis que la version scalée considérée est

\[
\mathcal L_R(\theta)
=
R\sum_{i=1}^n L_i(\theta),
\qquad
R=\frac2n.
\]

Ainsi,

\[
\mathcal L_R(\theta)=2\mathcal L(\theta),
\qquad
\arg\min_\theta\mathcal L_R(\theta)
=
\arg\min_\theta\mathcal L(\theta).
\]

Dans le code, en passant de $R=1/n$ à $R=2/n$, les quantités à modifier sont :

\begin{itemize}
    \item \textbf{Loss agrégée}
    \[
    \mathcal L_R
    =
    \frac{2}{n}\sum_i
    \frac{(y_i-Q_i)^2}{2\widetilde\sigma^2}
    =
    2\mathcal L.
    \]

    \item \textbf{Gradient de chaque couche}
    \[
    H_l
    =
    -\frac{R}{\widetilde\sigma^2}
    \sum_i (y_i-Q_i)\,
    r_{l,i}b_{l-1,i}^\top.
    \]
    Il est donc multiplié par $2$.

    \item \textbf{Facteur K-FAC d’activations}
    \[
    \widehat A_{l-1}
    =
    R\sum_i b_{l-1,i}b_{l-1,i}^\top
    =
    \frac{2}{n}\sum_i b_{l-1,i}b_{l-1,i}^\top.
    \]
    Il est multiplié par $2$.

    \item \textbf{Bloc K-FAC estimé}
    \[
    \widehat F_{ll}^{\mathrm{KFAC}}
    =
    \widehat A_{l-1}\otimes\widehat G_l.
    \]
    Il est donc multiplié par $2$.

    \item \textbf{Bloc exact estimé, si le code le calcule}
    \[
    \widehat F_{ll}
    =
    \frac{R}{\widetilde\sigma^2}
    \sum_i
    (b_{l-1,i}b_{l-1,i}^\top)
    \otimes
    (r_{l,i}r_{l,i}^\top).
    \]

    \item \textbf{Bloc exact de dernière couche}
    \[
    \widehat F_{LL}
    =
    \frac{R}{\widetilde\sigma^2}
    \sum_{a\in\mathcal A}
    \left(
    \sum_{i:a_i=a}
    b_{L-1,i}b_{L-1,i}^\top
    \right)
    \otimes e_ae_a^\top.
    \]

    \item \textbf{Empirical Fisher, s’il est implémenté}
    \[
    \widehat E
    =
    R\sum_i
    \nabla_\theta L_i\nabla_\theta L_i^\top.
    \]
    Il est multiplié par $2$.
\end{itemize}

Les quantités suivantes ne changent pas :

\begin{itemize}
    \item $D_v$, $g_l$ et $r_l$ ;

    \item $\widehat G_l$ :
    \[
    \widehat G_l
    =
    \frac{1}{n\widetilde\sigma^2}
    \sum_i r_{l,i}r_{l,i}^\top;
    \]

    \item $\widehat G_L=\frac{1}{n\widetilde\sigma^2}
    \operatorname{diag}(n_1,\ldots,n_{d_L})$ ;

    \item toutes les formules théoriques écrites avec des espérances.
\end{itemize}

Sans damping, $H_l$ et $\widehat A_{l-1}$ étant tous deux multipliés par $2$, la mise à jour

\[
-\eta\widehat G_l^\dagger H_l\widehat A_{l-1}^\dagger
\]

reste identique. Avec damping, ce changement d’échelle doit être répercuté dans le damping.
